%% Anmerkung zur Nutzung von Künstlicher Intelligenz (Leitlinie 4.6)

\chapter*{Anmerkung zur Nutzung von Künstlicher Intelligenz}

Generative KI-Werkzeuge wurden bei der Erstellung dieser Arbeit als Hilfsmittel
eingesetzt. Die folgende Übersicht legt offen, in welchen Arbeitsschritten dies
geschah und in welchem Umfang.

\vspace{1em}

\begin{tabularx}{\textwidth}{@{}p{3.2cm}XX@{}}
    \toprule
    \textbf{Werkzeug} & \textbf{Verwendungszweck} & \textbf{Umfang} \\
    \midrule
    Claude (Anthropic) & Überarbeitung eigener Textentwürfe, Vorschläge zur
    Gliederung und Kürzung &
    durchgehend, jeweils auf Grundlage eigener Entwürfe und nach eigener Prüfung
    übernommen \\
    \addlinespace
    Claude (Anthropic) & Formale Prüfung des \LaTeX-Quelltexts auf Verweise,
    Beschriftungen und Konsistenz & unterstützend \\
    \addlinespace
    Codex (OpenAI) & Kritische Prüfung der Argumentation und Quellenverwendung,
    ergänzende Literaturrecherche sowie Überarbeitung sachlicher Aussagen,
    insbesondere zu Prompting, Reproduzierbarkeit und Kosten & ausgewählte
    Abschnitte der Einleitung, Methodendarstellung, Auswertung und Diskussion \\
    \bottomrule
\end{tabularx}

\vspace{1.5em}

KI-Werkzeuge wurden auch zur kritischen Prüfung und Überarbeitung der Argumentation eingesetzt. Die berichteten empirischen Messwerte wurden dabei nicht durch KI erzeugt oder ersetzt, sie stammen aus den dokumentierten Versuchsläufen. Quellenhinweise und Formulierungsvorschläge sind anhand der Originalquellen und der eigenen Untersuchung zu prüfen. Die Verantwortung für die Quellenverwendung, die Interpretation der Ergebnisse und den gesamten Text liegt beim Verfasser.

Davon zu unterscheiden ist der Gegenstand dieser Arbeit selbst. Große
Sprachmodelle werden hier als Untersuchungsobjekt eingesetzt, um ihre Eignung
für \ac{ETL}-Aufgaben zu messen. Die dabei erzeugten Skripte und Ergebnistabellen
sind Messdaten und nicht Bestandteil des Textes.
